%=======================================================================
%  5. DISCUSSION & CONCLUSION  (text from the Word report; audit edits only)
%=======================================================================
\section{Discussion \& Conclusion}
\label{sec:discussion}

\subsection{Contribution of the biomimicry process}

Biomimicry helped us explore cooling functions before committing to familiar
building solutions. Studying how organisms manage heat encouraged us to abstract
and combine strategies, leading to the proposed link between façade shading, fluid
transport and a roof emitter (Sections~\ref{sec:discovering}--\ref{sec:creating}).
Without this process, we would probably have focused on improving existing
solutions individually.

\subsection{Strengths and weaknesses of the final concept}

V2 offers an unconventional approach for new buildings and especially retrofits:
shading is intended to limit façade heating, while the water loop transfers
collected heat to the roof. This aims to reduce heat release at street level
(Section~\ref{sec:morph}). However, circulation and the available roof area are the
main limitations. The first roof layout exceeded the available driving pressure and
was replaced by a parallel network (Section~\ref{sec:v2check}). Pump-free
circulation remains unverified, and refilling may require a pump.

\subsection{Assessment of the challenge and the KPIs}

The KPIs guided design decisions: the heat-flow estimates limited the connected
façade area (Section~\ref{sec:check2}). Against Table~\ref{tab:kpicompare}, KPI5 is
met on paper, KPI1 and KPI3 are estimates, and KPI2, KPI4 and KPI6 have no evidence
yet. Reflectance, shading duration, durability and actual cooling performance still
need verification.

\subsection{Innovative character and limits of the evidence}

The innovative aspect is the combination of biological principles in a maintainable
façade system intended for new buildings and retrofits. Our main lesson from the
course is that biological ideas become more useful when calculations and practical
constraints shape their technical transfer. Circulation and roof-output tests
(Table~\ref{tab:priorities}) are the priorities for assessing the concept's
potential.
